\section{Decomposition and inertia groups. \'Etale case}
\label{V.2}
Let $G$ be a
finite group acting on the right on the prescheme $X$. If
$x\in X$, one calls
% original p. 111
\emph{decomposition group of~$x$}
the stabilizer $G_d(x)$ of~$x$. This group acts canonically
(on the left) on the residue field
$\kres(x)$, and the set
of elements of~$G_d(x)$ that act trivially is
called the \emph{inertia group}
of~$x$, denoted~$G_i(x)$.

Suppose that $G$ acts on~$X$ in an admissible way and that
$Y$ is a prescheme over a prescheme
$Z$. Let us fix a $z\in Z$, and an algebraically
closed extension $\Omega$
of~$\kres(z)$ having a transcendence degree
greater than that of the $\kres(x)/\kres(z)$, where $x$ is a
point of~$X$ over~$z$. One can regard $\Spec(\Omega)$ as
a $Z$-scheme, and the points of~$X$ with values in $\Omega$
correspond to the homomorphisms
of~$\kres(z)$-algebras $\kres(x)\to\Omega$,
where $x$ is a point of~$X$ over~$z$;
since $\Omega$ has been taken large enough, every point $x$ of~$X$
over~$z$ is the locality of a point of~$X$ with values
in $\Omega$. If $X(\Omega)$ and $Y(\Omega)$ denote
respectively the set of points of~$X$ and $Y$ with values in
$\Omega$, one has a natural map
$$
X(\Omega)\to Y(\Omega),
$$
on the other hand, $G$ acts on~$X(\Omega)$ and the
preceding map is invariant under $G$. This being so, the
conclusions (ii) and (iii) of~\Ref{V.1.3} can also be interpreted
as follows: \emph{the preceding map is surjective and
identifies $Y(\Omega)$ with the quotient $X(\Omega)/G$. Moreover, if $x$ is
the locality of~$a\in X(\Omega)$, then the stabilizer of~$a$
in $G$ is none other than the inertia group $G_i(x)$}. All this is
moreover true without assuming $\Omega$ ``large enough''; this
last hypothesis serves only to ensure that one can
characterize the inertia group of every element of~$X$
over~$z$ as a ``geometric'' stabilizer. One
concludes from this, for instance, immediately:
\begin{proposition}
\label{V.2.1}
Let us make an extension of the base $Z'\to Z$, whence
$X'=X\times_ZZ'$. Let $x'$ be a point of~$X'$, $x$ its image in $X$;
then one has $G_i(x)=G_i(x')$.
\end{proposition}

It suffices, in the considerations above, to take for
$\Omega$ a sufficiently large extension
of~$\kres(z')$
(where $z,z'$ are the images of~$x,x'$ in $Z,Z'$).

\begin{proposition}
\label{V.2.2}
Under the conditions of~\Ref{V.1.3}, suppose $Y$ locally
noetherian, $X$ finite over~$Y$. Let $H$ be a subgroup of~$G$,
consider $X'=X/H$ (\cf \Ref{cor:V.1.7}), let
% original p. 112
$x\in X$, $x'$ its image in $X'$ and $y$ its image in $Y$.
\begin{enumerate}
\item[(i)] If $H\supset G_d(x)$, then the homomorphism
$\cal{O}_y\to\cal{O}_{x'}$ induces an isomorphism on the
completions.
\item[(ii)] If $H\supset G_i(x)$, then the homomorphism $\cal{O}_y\to
\cal{O}_{x'}$ is \'etale, \ie $X'$ is \'etale over~$Y$ at $x'$.
\end{enumerate}
\end{proposition}
Let $Y_1=\Spec(\widehat{\cal{O}_y})$; let us make the change of base
$Y_1\to Y$; one finds an $X_1=X\times_Y Y_1$ finite over~$Y_1$, on
which $G$ acts, the quotient being $Y_1$
by~\Ref{V.1.9}. Let~$y_1$ be the unique point of~$Y_1$ over~$y$;
since
$\kres(y)=\kres(y_1)$,
it follows that the fiber of~$X$ at~$y$
is isomorphic to that of~$X_1$ at $y_1$, whence a unique point
$x_1$ of~$X_1$ over~$x$. Moreover, by~\Ref{V.1.9} one will have
$X_1/H=X'_1=X'\times_Y Y_1$; let $x'_1$ be the image of~$x_1$ in $X'_1$;
it lies over~$x'$, and one easily checks ($X'$ being
of finite type over~$Y$) that the homomorphism
$\cal{O}_{x'}\to\cal{O}_{x'_1}$ induces an isomorphism on the
completions. Hence one is reduced to the case where $Y$ is the
spectrum of a complete local ring, say $B$, hence~$X$ the spectrum of a
ring $A$ finite over~$B$, composed of a finite number of local rings~$A_x$ corresponding
to the points $x_i$ of~$X$ over~$Y$. If $A_0$ corresponds to $x=x_0$,
then $A$ is identified with the ring $\Hom_{G_d}(G,A_0)$ of the functions
$f\colon G\to A_0$ such that $f(st)=sf(t)$ for $s\in G_d$, the
operations of~$G$ on these functions being defined by
$(uf)(t)=f(tu)$. One sees therefore that if $H$ is an arbitrary subgroup
of~$G$, then $A^H$ is the ring of the functions $f\colon G\to A_0$
such that
$$
f(stu)=sf(t)\qquad\text{for}\ s\in G_d,\ u\in H
$$
hence it is a semi-local ring whose local components
correspond to the double cosets $G_daH$ in $G$, the double
coset defined by $a\in G$ corresponding (thanks to
the map $f\mto f(a)$) to the subring $A_0^{H(a)}$ of~$A_0$,
where $H(a)=G_d\cap aHa^{-1}$. Moreover, the local component of
$A^H$ corresponding to the image $x'$ of~$x$ is also the one
corresponding to the double coset $G_dH$ of the neutral
element; its local component is therefore $A_0^{G_d\cap H}$. If therefore
$G_d\subset H$,
one finds $A_0^{G_d}=A^G=B$, which proves (i). To
prove (ii), one can, by passing to a suitable finite flat
extension of~$A$, and using~\Ref{V.2.1}, reduce to the case where
the residue extension
$\kres(x)/\kres(y)$
is trivial. But
then $G_i(x)=G_d(x)$, and one is reduced to the preceding case.

\begin{corollary}
\label{V.2.3}
Under
% original p. 113
the conditions of~\Ref{V.2.2}, suppose $G_i(x)=(e)$; then X
is \'etale over~$Y$ at $x$. Hence if $G_i(x)=(e)$ for every $x \in
X$, then $X \to Y$ is an \'etale morphism.
\end{corollary}

There is a partial converse:

\begin{corollary}
\label{V.2.4} Suppose $X$ connected and the group $G$ faithful
on~$X$. For $p \colon X \to Y=X/G$ to be \'etale, it is necessary and
sufficient that the inertia groups of the points of~$X$ be reduced
to the neutral element. If this is so, $G$ is identified with the
group of all $Y$-automorphisms of the $Y$-scheme~$X$.
\end{corollary}

Taking~\Ref{V.2.3} into account, one can suppose $X$ \'etale over~$Y$. But if an $s \in G$ is in some $G_i(x)$, it then follows
from I~\Ref{I.5.4} that $s$ acts trivially on~$G$, hence is
the unit element since $G$ is faithful, which proves
the first assertion. Let $u$ be a $Y$-automorphism of~$X$, let
$x \in X$. By proposition~\Ref{V.1.3}, there exists an $s \in
G$ such that $s(x)=u(x)$, and inducing the same residue
homomorphism
$\kres(x)/\kres(x')$
as $u$. By \loccit one therefore has
$s=u$, which completes the proof.

\begin{remark}
\label{V.2.5}
The hypothesis that $G$ acts faithfully is obviously
not superfluous in corollary~\Ref{V.2.4}. The same holds
for the hypothesis that $X$ is connected, as one sees for instance by
taking $X=Y \times E$, $E$ being a finite set, and $G$ the
group of permutations of~$E$: $G$ acts with plenty of inertia,
nevertheless $(Y \times E)/G=Y \times (E/G)=Y$, and $X$ is \'etale
over~$Y$. Taking for $G$ a group strictly smaller than the
symmetric group of~$E$, but acting transitively on~$E$,
one sees that there will also be $Y$-automorphisms of~$X$ not coming
from~$G$.
\end{remark}

The typical example of a group $G$ acting without inertia is that of
$Y \times G$, on which one makes $G$ act by means of its
operations on the factor $G$ by right translations: a
$Y$-prescheme $X$ with right group of operators
$G$ is said to be \emph{trivial} if it is isomorphic to $Y \times G$.

To make the link between preschemes with finite groups
of operators and the notion of principal bundle in a
category (a link which, moreover, we shall not need for the
rest of the seminar, but which is important in other contexts) the
following considerations are useful. We fix a
base prescheme $Y$, and we place ourselves in the
category of $Y$-preschemes. If $G$ is a finite
% original p. 114
group, we shall set for brevity $G_Y=Y \times G$; it is thus a
finite group scheme over~$Y$ (\cf \numero 1), and if $X$ is
a $Y$-prescheme, one has
$$
X \times_Y G_Y = X \times G
$$
(same reference). The datum of a $Y$-morphism $X
\times_Y G_Y \to X$ is therefore equivalent to the datum of a
$Y$-morphism $X \times G \to X$, \ie to the datum, for every
$g \in G$, of a $Y$-morphism $T_g \colon X \to X$. One notes
immediately that for the datum of the $T_g$ to define on~$X$ a structure of prescheme with right group of operators
$G$ (\ie $T_{gg'}=T_{g'}T_g$, $T_e=\id_x$) it is necessary and
sufficient that the corresponding $Y$-morphism $X \times_Y G_Y \to X$
define on~$X$ a structure of~$Y$-prescheme with
$Y$-group scheme of operators (in the general sense
of objects with $\cal{C}$-group of operators in a
category $\cal{C}$). Suppose that this is so. Let us recall that
$X$ is said to be \emph{formally principal homogeneous}
under~$G_Y$\kern1pt\footnote{one now says rather: $X$ is a
pseudo-torsor under $G_Y$.} if the canonical morphism
$$
X \times_Y G_Y \to X \times_Y X
$$
whose components are respectively $\pr_1$ and the
multiplication morphism $\pi \colon X \times_Y G_Y \to X$, is an
isomorphism. In the present case, identifying the left-hand side with $X
\times G$, the morphism considered is the one which, to every $g
\in G$, associates the morphism
$$
(\id_X,T_g)=(\id_X \times_Y T_g)\Delta_{X/Y} \colon X \to X \times_Y X
$$
and consequently, to say that $X$ is formally principal homogeneous
under $G_Y$ also means that $X \times_Y X$ is isomorphic to the
direct sum of the transforms of the diagonal by the
elements $(e,g)$ of~$G \times G$ (acting on~$X \times_Y
X$ in the obvious way), where $e$ denotes the unit
element of~$G$. If one wishes not to distinguish left and right,
and to give a formula that remains applicable to the product of more than two
factors identical to $X$, one can formulate the condition by saying
that the canonical morphism
$$
X \times_G (G \times G) \to X \times_Y X
$$
obtained by attaching to the pair $(g,g')$ the morphism
$$
(T_g,T_{g'})=(T_g \times_Y T_{g'})\Delta_{X/Y} \colon X \to X \times_Y
X
$$
and by making $G$ act on the left on~$G \times G$ via
the diagonal homomorphism;
$$
s(g,g')=(sg,sg'),
$$
is an \emph{isomorphism}.

The
% original p. 115
notion of \emph{principal homogeneous space}
is deduced from that of formally principal
homogeneous space by adding a supplementary axiom, ensuring that the
``quotient'' of~$X$ by $G_Y$ exists and is precisely
the right unit object of the category, here $Y$. This
axiom may vary according to the context, and is often made explicit most
conveniently (in the yoga of ``descent'') by requiring that
the object with operators become ``trivial'', \ie isomorphic to the
product $X \times_Y G_Y$ (in the present case $X \times G$) by a
suitable change of base, of a specified type (in such a way, in practice, as to permit the technique of descent;
\cf Grothendieck, Technique de descente et th\'eor\`emes
d'existence en G\'eom\'etrie Alg\'ebrique, S\'em. Bourbaki
\No 190, pages 26 to 28)\footnote{\Cf Exp\ptbl \Ref{VIII} for the
theory of flat descent.}. In this order of ideas,
let us point out here the characterization of principal
homogeneous bundles with group $G$ (in the sense of \loccit):

\begin{proposition}
\label{V.2.6}
Let $Y$ be a locally noetherian prescheme, $X$ a
$Y$-prescheme with finite group $G$ of operators
acting on the right. The following conditions are
equivalent:
\begin{enumerate}
\item[(i)] $X$ is finite over~$Y$, $Y=X/G$, the inertia groups of the
points of~$X$ are reduced to the unit element.
\item[(ii)] There exists a faithfully flat and
quasi-compact change of base $Y_1 \to Y$ such that $X_1=X \times_Y Y_1$ is a
trivial $Y_1$-prescheme with operators, \ie isomorphic
to $Y_1 \times G$.
\item[(ii~bis)] Like (ii), but with $Y_1 \to Y$ being finite,
\'etale, surjective.
\item[(iii)] $X$ is formally principal homogeneous under $G_Y$,
and faithfully flat and quasi-compact over~$Y$.
\end{enumerate}
\end{proposition}
\textit{Proof} (i)$\To$(ii~bis) We shall take
$Y_1=X$, noting that $X \to Y$ is indeed finite, \'etale by~\Ref{V.2.3},
and surjective. Let us show that $X_1$ is then trivial over~$Y_1$, which
will follow from the

\begin{corollary}
\label{V.2.7} If \textup{(i)} is satisfied and if $X$ admits a
section over~$Y$, then $X$ is a trivial space with operators.
\end{corollary}

Indeed, this section makes it possible to define a $G$-morphism $X
\times G \to X$, surjective since $G$ is transitive on the fibers of
$X$, injective since $G$ acts without inertia; finally, it is a
local isomorphism by virtue of I~\Ref{I.5.3} since $X$ is \'etale
over~$Y$. Hence it is an isomorphism.

(ii~bis)
% original p. 116
trivially implies (ii), which implies (i) because the ingredients
of (i) are ``invariant'' under faithfully flat
quasi-compact extension of the base (\cf the S\'eminaire Bourbaki cited
above for ``finite''; for the inertia groups, one
applies~\Ref{V.2.1}, and for $Y=X/G$, a converse
of~\Ref{V.1.9} in the case of a \emph{faithfully flat} change of base,
which we had forgotten to make explicit).

We have proved (i)$\To$(iii) in passing, while proving
(i)$\To$(ii~bis). Finally (iii)$\To$(ii), because the
first hypothesis in (iii) means precisely that
$X$ becomes trivial on making the change of base $Y_1=X$; whence
(ii) since $X$ is faithfully flat and quasi-compact over~$Y$.

\begin{definition}
\label{V.2.8} A $Y$-prescheme $X$ with right group
of operators $G$ satisfying the
equivalent conditions of~\Ref{V.2.6} is called a \emph{principal
covering of~$Y$, with Galois group~$G$}.
\end{definition}

\section{Automorphisms and morphisms of \'etale coverings}
\label{V.3}

\begin{proposition}
\label{V.3.1} Let $X$ be \'etale, separated, of finite type over~$Y$, with $Y$ locally noetherian, and let $G$ be a finite group acting on~$X$ by $Y$-automorphisms. Then $G$ acts in an
admissible way and the quotient prescheme $X/G$ is \'etale over~$Y$.
\end{proposition}

One does not suppose $X$ finite over~$Y$; however $X$ is quasi-projective
over~$Y$, whence the existence of~$X/G$ thanks
to~\Ref{V.1.8}. Let us first prove

\begin{corollary}
\label{V.3.2} The morphism $X \to X/G$ is \'etale.
\end{corollary}

One can obviously suppose $G$ transitive on the set of
connected components of~$X$, then, by considering the
stabilizer of a connected component, that $X$ is even
connected. Finally, one can suppose that $G$ acts faithfully.
But then one sees as in~\Ref{V.2.4} that $G$ acts without
inertia, hence by~\Ref{V.2.3} it follows that $X \to X/G$ is
\'etale. One concludes thanks to the

\begin{lemma}[afterthought to Expos\'e I]
\label{V.3.3} Let $X \to X' \to
Y$ be morphisms of finite type, $x$ a point of~$X$, $x'$ and $y$ its
images. Suppose $Y$ locally noetherian. If two of the
morphisms considered are \'etale at the marked points, so is
the third.
\end{lemma}

It remains only to look at the case where $X \to X'$ and $X \to
Y$ are \'etale at $x$, and to prove that $X' \to Y$ is \'etale at $x'$ (which
is the case that we need
% original p. 117
for~\Ref{V.3.1}). Making a suitable flat extension of the base
$Y$, one is reduced to the case where the residue extension
$\kres(x)/\kres(y)$
is trivial. Consider the homomorphisms
$\cal{O}_y \to \cal{O}_{x'} \to \cal{O}_x$ and the homomorphisms
deduced by passing to the completions; the hypothesis
means that $\hat{\cal{O}_y} \to \hat{\cal{O}_x}$ and
$\hat{\cal{O}_{x'}} \to \hat{\cal{O}_x}$ are isomorphisms,
whence immediately that $\hat{\cal{O}_y} \to \hat{\cal{O}_{x'}}$
is one, which proves the lemma.

\begin{corollary}
\label{V.3.4} If $X$ is finite and \'etale over~$Y$, then $X/G$
is finite and \'etale over~$Y$.
\end{corollary}

\begin{proposition}
\label{V.3.5} Let $X$, $X'$ be two \'etale coverings of
$Y$. Then every $Y$-morphism $f \colon X \to X'$ factors as the
product of a surjective \'etale morphism $X \to X''$ and
the canonical immersion $X'' \to X'$ of a subset $X''$ of~$X'$ that is
both open and closed.
\end{proposition}

One knows (I~\Ref{I.4.8}) that $f$ is \'etale, hence an open
morphism; on the other hand, $X$ being finite over~$Y$, $f$ is closed,
hence $f(X)=X''$ is a subset of~$X'$ that is both open and closed. We are done (N.B. it would have sufficed that $X'$, instead of being an
\'etale covering, be unramified over~$Y$).

\begin{corollary}
\label{V.3.6} With the preceding notation, $X \to X'$
is a strict epimorphism in the category of
preschemes, and $X' \to X''$ is a monomorphism (and even
a strict monomorphism) in the category of preschemes.
\end{corollary}

The first assertion means by definition that the sequence of
morphisms
$$
\xymatrix@C=.7cm{X \times_{X''} X \ar@<2pt>[r]^-{\pr_1}
\ar@<-2pt>[r]_-{\pr_2} & X \ar[r] & X''}
$$
is exact, and this follows from the fact that $X \to X''$ is finite and
faithfully flat, as one sees easily (\cf Grothendieck, \loccit). The dual assertion for $X'' \to X'$ is even more trivial.

Corollary~\Ref{V.3.6} will be useful to us for the theory of the
fundamental group in the next no.; it is possible (for those
who do not like the notion of strict epimorphism) to replace
corollary~\Ref{V.3.6} by such a variant as the reader will arrange
to his personal taste. Let us merely take the opportunity to
point out that a factorization
% original p. 118
$f=f'f''$, with $f''$ a strict epimorphism and $f'$ a
monomorphism, is necessarily unique up to unique isomorphism
(in any category); however, there may exist at the
same time a factorization $f=f_1f_2$ having the
dual properties: $f_2$ is an epimorphism, $f_1$ a
strict monomorphism (also unique up to unique
isomorphism), which is not isomorphic to the preceding one: it
suffices to take for instance the category of topological vector spaces
(Hausdorff, if one insists), and for $u \colon X \to
X'$ a morphism such that $u(X)$ is not closed.

\begin{proposition}
\label{V.3.7} Let $Y$ be a \emph{connected} locally
noetherian prescheme, $y$ a point of~$Y$, $\mathit{\Omega}$ an
algebraically closed extension
of~$\kres(y)$. For every $X$ over~$Y$, one
denotes by $X(\Omega)$ the set of points of~$X$
with values in $\Omega$. Let $X$, $X'$ be \'etale coverings
of~$Y$ and $u \colon X \to X'$ a $Y$-morphism such that
the corresponding map $X(\Omega) \to
X'(\Omega)$ is an isomorphism. Then $u$ is an
isomorphism.
\end{proposition}

One is immediately reduced to the case where $X'$ is
connected. Since $X \to X'$ is finite and \'etale, one knows that the
geometric number of points in a fiber of~$X \to X'$ is
constant, and equal to 1 if and only if the morphism
considered is an isomorphism. Now this number is also the
number of elements in a fiber of~$X(\Omega) \to
X'(\Omega)$, whence the conclusion.
